% Generated draft: requires mathematical and source review.
\documentclass[11pt]{article}
\usepackage[margin=25mm]{geometry}
\usepackage{fontspec}
\setmainfont{Latin Modern Roman}
\usepackage{amsmath,amssymb,mathtools,unicode-math}
\setmathfont{Latin Modern Math}
\usepackage{enumitem,needspace,xurl,hyperref}
\setlength{\parindent}{0pt}
\setlength{\parskip}{5pt}
\setlength{\emergencystretch}{3em}
\newcommand{\N}{\mathbb{N}}
\newcommand{\Z}{\mathbb{Z}}
\newcommand{\Q}{\mathbb{Q}}
\newcommand{\R}{\mathbb{R}}
\newcommand{\C}{\mathbb{C}}
\newcommand{\F}{\mathbb{F}}
\newcommand{\A}{\mathbb{A}}
\newcommand{\PP}{\mathbb{P}}
\newcommand{\E}{\mathbb{E}}
\newcommand{\eps}{\varepsilon}
\newcommand{\dd}{\,\mathrm{d}}
\begin{document}
\section*{UnsolvedMath ID: 6200108}
\textbf{Source problem code:} catalogue entry
\subsection{Definitions and mathematical statement}
A finitely generated group $G$ is word hyperbolic if its Cayley graph, with the path metric from a finite generating set, has uniformly thin geodesic triangles.

The supplied item asks whether the Coarse Whitehead Conjecture holds under this hyperbolicity hypothesis on $G$. The conjecture itself is not stated in the record. In particular no complex associated to $G$, coarse homotopy theory, map or equivalence category, scale quantifiers, or conclusion is supplied. Consequently a standalone implication expressing the conjecture's hyperbolic case cannot be recovered from these data.

The research summary sketches a finitely presented group with ``coarsely trivial'' fundamental group of an unspecified complex and trivial higher coarse homotopy, followed by a coarse-triviality conclusion. That sketch does not define the conjecture and is not promoted to an exact theorem statement. The target retained here is to establish or refute the original Coarse Whitehead Conjecture when $G$ is hyperbolic, once its exact source formulation is supplied.
\subsection{Short English statement}
Does the Coarse Whitehead Conjecture hold for hyperbolic groups? The required named conjecture is missing, so the record does not support a complete quantified formulation. The supplied summary calls the hyperbolic case unresolved but gives only an incomplete paraphrase. Its OPEN-TRIAGE label and the 2005 source are retained with current status marked for review.
\subsection{Sources}
\begin{itemize}
\item Local UnsolvedMath registry, permanent ID 6200108, source code catalogue entry, ``Boundaries of Groups and Kleinian Groups --- Problem 108,'' proposed year 2005; no proposer is specified. Dataset curation and metadata: Ulam AI and UnsolvedMath Contributors, CC BY 4.0.
\item Kapovich, \emph{Problems on Boundaries of Groups and Kleinian Groups} (2005), Problem 108, PDF page 27, \url{https://www.math.ucdavis.edu/~kapovich/EPR/problems.pdf}; normalized pdf-text extraction, accessed 2026-07-29.
\item The research summary attributes the conjecture to Kapovich/coarse group geometry but gives no original statement or separate reference. Source-specific provenance marks current status and statement-reuse terms NEEDS\_REVIEW for the public author-hosted PDF, separately from the dataset licence.
\end{itemize}
\end{document}
